\documentclass[12pt,a4paper]{article}
\usepackage[margin=1in]{geometry}
\usepackage{xcolor}
\usepackage{enumitem}
\usepackage{hyperref}
\usepackage{titlesec}

% Configure hyperlinks
\hypersetup{
    colorlinks=true,
    linkcolor=blue,
    filecolor=magenta,      
    urlcolor=cyan,
    pdftitle={Frontend Folder Structure},
}

% Title formatting
\titleformat{\section}
{\normalfont\Large\bfseries\color{blue!70!black}}{\thesection}{1em}{}
\titleformat{\subsection}
{\normalfont\large\bfseries\color{blue!50!black}}{\thesubsection}{1em}{}

% Document Information
\title{\textbf{Frontend Folder Structure} \\ \Large Daily Utility Tools Application}
\author{Flutter Application}
\date{\today}

\begin{document}

\maketitle
\tableofcontents
\newpage

\section{Overview}
This document describes the high-level folder structure of the Frontend Flutter application for the Daily Utility Tools project.

\section{Root Level Folders}

\subsection{lib/}
\textbf{Purpose:} Contains all the Dart application source code.

This is the main application folder where all the Flutter/Dart code lives.

\subsubsection{lib/api\_services/}
Handles communication with the backend API. Contains API route definitions and service classes for each feature.

\subsubsection{lib/cards/}
Reusable card components for displaying results. Each feature has its own subfolder with result display cards.

\subsubsection{lib/models/}
Data models and data transfer objects (DTOs) that represent the structure of data used in the app.

\subsubsection{lib/notifiers/}
State management logic using ChangeNotifier pattern. Each feature has its own notifier to manage state.

\subsubsection{lib/providers/}
Provider setup for dependency injection. Links notifiers to the widget tree.

\subsubsection{lib/routes/}
Navigation and routing configuration. Defines app routes and router setup.

\subsubsection{lib/screens/}
Main screens/pages of the application. Each utility tool has its own screen.

\subsubsection{lib/theme/}
Application theming configuration including colors, fonts, and styling.

\subsubsection{lib/widgets/}
Reusable UI components organized by feature. Each feature has its own subfolder with custom widgets.

\subsection{assets/}
\textbf{Purpose:} Static resources like images, fonts, and other files.

Contains all non-code resources that the app uses.

\subsubsection{assets/images/}
Image files used throughout the application (icons, backgrounds, illustrations).

\subsection{test/}
\textbf{Purpose:} Testing code including unit tests and widget tests.

Contains all test files for the application.

\subsection{android/}
\textbf{Purpose:} Android platform-specific configuration and native code.

Contains Gradle build files, Android manifest, and native Android code needed for the app to run on Android devices.

\subsection{ios/}
\textbf{Purpose:} iOS platform-specific configuration and native code.

Contains Xcode project files, Info.plist, and native iOS code needed for the app to run on iOS devices.

\subsection{web/}
\textbf{Purpose:} Web platform-specific configuration.

Contains HTML, manifest, and web assets for running the app as a Progressive Web App (PWA).

\subsection{windows/}
\textbf{Purpose:} Windows platform-specific configuration and native code.

Contains CMake configuration and C++ code for running the app on Windows desktop.

\subsection{linux/}
\textbf{Purpose:} Linux platform-specific configuration and native code.

Contains CMake configuration and C++ code for running the app on Linux desktop.

\subsection{macos/}
\textbf{Purpose:} macOS platform-specific configuration and native code.

Contains Xcode project files and Swift code for running the app on macOS.

\section{Application Architecture}

\subsection{How the Folders Work Together}

\begin{enumerate}
    \item \textbf{screens/} contain the main UI pages users see
    \item \textbf{widgets/} contain reusable UI components used in screens
    \item \textbf{cards/} contain specialized components for displaying results
    \item \textbf{notifiers/} manage the state and business logic
    \item \textbf{providers/} connect notifiers to the UI
    \item \textbf{api\_services/} handle backend communication
    \item \textbf{models/} define data structures
    \item \textbf{routes/} handle navigation between screens
    \item \textbf{theme/} provides consistent styling
\end{enumerate}

\subsection{Feature Organization}

Each utility tool (Age Calculator, EMI Calculator, Barcode Generator, QR Generator, Unit Converter) follows the same pattern:

\begin{itemize}
    \item One screen in \texttt{screens/}
    \item Custom widgets in \texttt{widgets/[feature\_name]/}
    \item Result cards in \texttt{cards/[feature\_name]/}
    \item One model in \texttt{models/}
    \item One notifier in \texttt{notifiers/}
    \item One provider in \texttt{providers/}
    \item One service in \texttt{api\_services/services/}
\end{itemize}

\section{Platform Folders Purpose}

All platform folders (\texttt{android/}, \texttt{ios/}, \texttt{web/}, \texttt{windows/}, \texttt{linux/}, \texttt{macos/}) serve the same purpose: they contain platform-specific code and configuration that Flutter needs to run the app on that particular platform. You generally don't modify these unless you're adding native features.

\section{Summary}

\begin{itemize}
    \item \textbf{lib/} = Your app code
    \item \textbf{assets/} = Images and resources
    \item \textbf{test/} = Tests
    \item \textbf{Platform folders} = Configuration for each OS
\end{itemize}

\end{document}
